% SGA 3, Expose VI_A, English reconstruction -- Loop 1.
% Scope: component-local pages 1--4, from the title through the end of
% Corollary 0.5.2, stopping before the Section 1 heading.
% Combined-reader pages: 304--307. Printed pages: 293--296.
% Authority: Polo--Gille Exp6A-13oct24.pdf,
% SHA-256 6C8795572D1FC3B21BD02D24B1BFB20A8FB447475702C2674C96939B1EB69582.
% Jacob Reinhold's CC BY 4.0 English Markdown was consulted only as a
% comparison lineage. The Polo--Gille PDF controls prose, formulas,
% numbering, footnotes, and the source-bound Loop-1 diagram.

\setcounter{footnote}{-1}
\chapter*{第 \texorpdfstring{VI\textsubscript{A}}{VI_A} 篇\\
代数群的一般理论}
\phantomsection
\label{sga3:VIA:title}
\addcontentsline{toc}{chapter}{第 VI-A 篇. 代数群的一般理论}
\markboth{第 VI\textsubscript{A} 篇. 代数群的一般理论}%
  {第 VI\textsubscript{A} 篇. 代数群的一般理论}

\begin{center}
  \textit{P. Gabriel 著}
\end{center}

% Authority: Exp6A local p. 1; combined-reader p. 304; printed p. 293.
\footnote{编者注: 2024 年 10 月 13 日版本.}
本章中, \(A\) 都表示剩余域为 \(k\) 的 Artin 局部环.
\(\operatorname{Spec} A\) 上的群概形 \(G\) 简称为
\(A\)-\emph{群}. 若其底概形在 \(A\) 上局部有限型,
则称此 \(A\)-群\emph{局部有限型};
若其底概形在 \(A\) 上有限型, 则称其为\emph{代数的}.

\section*{0. 预备说明}
\phantomsection
\label{sga3:VIA:0}
\addcontentsline{toc}{section}{0. 预备说明}

\medskip
\noindent\textbf{0.1.}
\phantomsection
\label{sga3:VIA:0.1}
先考虑任意概形 \(S\) 上的一个群概形 \(G\).
称结构态射
\[
  \mu:G\times_S G\longrightarrow G
\]
为\emph{乘法}; 称态射 \(c:G\to G\) 为\emph{求逆},
其中对 \(S\) 上的每个概形 \(T\) 及每个 \(x\in G(T)\),
有 \(c(T)(x)=x^{-1}\).
若 \(U\) 与 \(V\) 是 \(G\) 的底集合的子集,
以 \(U\cdot V\) 表示 \(G\times_S G\) 的如下子集在乘法下的像:
其第一投影属于 \(U\), 第二投影属于 \(V\) 的那些点所成的子集.
同样地, 记号 \(U^{-1}\) 与 \(c(U)\) 等价.

设 \(\operatorname{pr}_1\) 是 \(G\times_S G\) 到第一因子的投影,
设
\[
  \sigma:G\times_S G\longrightarrow G\times_S G
\]
是以 \(\operatorname{pr}_1\) 和 \(\mu\) 为分量的态射.
对每个 \(S\)-概形 \(T\), 映射 \(\sigma(T)\) 为
\[
  (x,y)\longmapsto (x,xy);
\]
所以 \(\sigma\) 是自同构.
此自同构与 \(G\times_S G\) 到第二因子的投影
\(\operatorname{pr}_2\) 的复合就是乘法态射.
当 \(G\) 在 \(S\) 上平坦时, \(\operatorname{pr}_2\),
因而 \(\mu\), 都是平坦态射;
当 \(G\) 在 \(S\) 上平滑时, \(\operatorname{pr}_2\),
因而 \(\mu\), 都是平滑态射, 其余类推.

\medskip
\noindent\textbf{0.2.}
\phantomsection
\label{sga3:VIA:0.2}
从现在起, 假设 \(S\) 是剩余域为 \(k\) 的 Artin 局部环 \(A\) 的谱.
以 \((\mathbf{Sch}_{/k})_{\mathrm{red}}\) 表示 \(k\) 上的约化概形范畴.
对 \(A\) 上的每个概形 \(X\), 约化概形 \(X_{\mathrm{red}}\)
是 \((\mathbf{Sch}_{/k})_{\mathrm{red}}\) 的对象,
且函子
\[
  X\longmapsto X_{\mathrm{red}}
\]
是 \((\mathbf{Sch}_{/k})_{\mathrm{red}}\) 到
\(\mathbf{Sch}_{/A}\) 的包含函子的右伴随.
所以, 对每个 \(A\)-群 \(G\), \(G_{\mathrm{red}}\)
是范畴 \((\mathbf{Sch}_{/k})_{\mathrm{red}}\) 中的一个群.
\footnote{编者注: 以下几句为新增内容.}
换言之, 对每个约化 \(k\)-概形 \(T\),
\(G_{\mathrm{red}}(T)\) 配备了一个关于 \(T\) 函子性的群结构.
必须记住, \(G_{\mathrm{red}}\) 未必是一个 \(k\)-群,
因为其 ``乘法'' 只是一个态射
\stepcounter{footnote}%
\[
  (G_{\mathrm{red}}\times_k G_{\mathrm{red}})_{\mathrm{red}}
  \longrightarrow
  G_{\mathrm{red}}\text{\hyperlink{sga3:VIA:fn:red-examples}{%
    \textsuperscript{\normalfont\thefootnote}}}.
\]
\footnotetext[\value{footnote}]{\hypertarget{sga3:VIA:fn:red-examples}{}编者注:
非完美域 \(k\) 上使得 \(G_{\mathrm{red}}\) 不是 \(k\) 上的群概形的
群概形 \(G\) 的例子, 见下文 \SGARefLink{ex:sga3-VIA-1.3.2}{1.3.2}.}

\footnote{编者注: \SGARefLink{sga3:VIA:0.2}{0.2} 的余下部分及
\SGARefLink{sga3:VIA:0.3}{0.3} 作了少量修改.}
然而, 若 \(k\) 是完美域, 则
\((\mathbf{Sch}_{/k})_{\mathrm{red}}\) 到 \(\mathbf{Sch}_{/k}\)
的包含函子与积可交换.
所以, \((\mathbf{Sch}_{/k})_{\mathrm{red}}\) 中的群可等同于
\(k\) 上底概形约化的群概形.
此时, 若 \(G\) 是一个 \(k\)-群,
则 \(G_{\mathrm{red}}\) 是 \(G\) 的一个群子概形;
一般来说, 此子群在 \(G\) 中并不正规.

例如, 若 \(k\) 是特征为 \(3\) 的域,
则常群 \((\mathbb Z/2\mathbb Z)_k\) 在可对角化群
\(D_k(\mathbb Z/3\mathbb Z)\) 上有非平凡作用
(参见 \SGARefLink{sga3:i:subsection:4-1}{第 I 篇, 4.1} 和
\SGARefLink{cor:sga3-VIB-4.4}{4.4}).
若以 \(G\) 表示由此作用定义的 \(D_k(\mathbb Z/3\mathbb Z)\)
与 \((\mathbb Z/2\mathbb Z)_k\) 的半直积,
则 \(G_{\mathrm{red}}\) 等同于 \((\mathbb Z/2\mathbb Z)_k\),
且在 \(G\) 中不正规.
\footnote{编者注: 若要 \(G\) 连通的类似例子,
可在 \(\operatorname{char}(k)=p>0\) 时,
取 \(\alpha_{p,k}\) 与 \(\mathbb G_{m,k}\) 的半直积.}

设 \(k\) 是任意域, \(k^{p^{-\infty}}\) 是其完美闭包,
\(H\) 是 \((\mathbf{Sch}_{/k})_{\mathrm{red}}\) 中的一个群. 则
\[
  (H\otimes_k k^{p^{-\infty}})_{\mathrm{red}}
\]
是完美域 \(k^{p^{-\infty}}\) 上的一个群概形.
由于 \(H\otimes_k k^{p^{-\infty}}\) 与
\((H\otimes_k k^{p^{-\infty}})_{\mathrm{red}}\)
有相同的底拓扑空间,
\((\mathbf{Sch}_{/k})_{\mathrm{red}}\) 中的群与 \(k\)-群
共有某些在基域扩张下不变的拓扑性质.
例如, 由 \hyperref[sga3:VIA:0.3]{0.3} 及以上说明将可推出,
\((\mathbf{Sch}_{/k})_{\mathrm{red}}\) 中的每个群都是分离的.

以下, 当我们处理一个 \(A\)-群 \(G\) 的非空局部闭子集 \(U\),
且它满足
\[
  U\cdot U=U
  \qquad\text{及}\qquad
  U^{-1}=U.
\]
时, 就会遇到 \((\mathbf{Sch}_{/k})_{\mathrm{red}}\) 中的群.
事实上, \(G\) 的由 \(U\) 定义的约化子概形,
此时就是 \((\mathbf{Sch}_{/k})_{\mathrm{red}}\) 中的一个群.

% Authority: Exp6A local p. 2; combined-reader p. 305; printed p. 294.
\medskip
\noindent\textbf{0.3.}
\phantomsection
\label{sga3:VIA:0.3}
\textit{\(A\)-群 \(G\) 总是分离的,
因为单位截面 \(e:\operatorname{Spec}A\to G\) 是闭浸入.}
事实上, 设 \(x\) 是 \(\operatorname{Spec}A\) 的唯一点,
\(\eta:G\to\operatorname{Spec}A\) 是结构态射.
由于 \(\eta\circ e=\operatorname{id}_{\operatorname{Spec}A}\),
对 \(G\) 的每个包含 \(e(x)\) 的仿射开集 \(U=\operatorname{Spec}B\),
同态 \(B\to A\) 都有一个截面, 从而是满射.
因此 \(e\) 是闭浸入.
\footnote{编者注: 这个论证对每个零维局部环 \(A\) 都有效,
并表明当 \(S\) 是离散概形时, 每个 \(S\)-群都是分离的;
参见 VI\textsubscript{B}, \SGARefLink{prop:sga3-VIB-5.2}{5.2}.
另一方面 (参见 VI\textsubscript{B}, \SGARefLink{ctr:sga3-VIB-5.6.4}{5.6.4}),
若 \(S\) 含有一个非孤立闭点 \(s\),
则将 \(S\) 的两个副本沿开子集 \(S-\{s\}\) 粘合所得的
\(S\)-概形 \(G\) 在 \(S\) 上并不分离,
尽管它带有一个 \(S\)-群结构.}

\(G\times_A G\) 中的对角等同于从
\((\mathbf{Sch}_{/A})^\circ\) 到 \(\mathbf{Sets}\) 的如下函子:
它将每个 \(A\)-概形 \(S\) 映为 \(G(S)\) 的单位元在映射
\[
  \varphi(S):G(S)\times G(S)\longrightarrow G(S),
  \qquad (x,y)\longmapsto x\cdot y^{-1}.
\]
下的逆像. 因此有下列笛卡尔方块.
由于对角态射由一个闭浸入经基变换得到, 它本身也是闭浸入.

\begin{center}
  \phantomsection
  \label{sga3:VIA:0.3-cartesian-square}
  \begin{tikzcd}[row sep=large,column sep=large]
    G\times_S G
      \arrow[r,"\varphi"]
    &
    G
    \\
    G
      \arrow[u,"\text{对角态射}"]
      \arrow[r]
    &
    \operatorname{Spec}A
      \arrow[u,"\text{单位截面}"']
  \end{tikzcd}
\end{center}

% The exact source-PDF crop above prints G x_S G in its upper-left corner.
% The surrounding source prose identifies the fiber product as G x_A G.

% Authority: Exp6A local p. 3; combined-reader p. 306; printed p. 295.
\medskip
\noindent\textbf{0.4.}\footnote{编者注: 以下内容对原文作了扩充;
特别新增了\hyperref[sga3:VIA:0.4.1]{注 0.4.1}.}
\phantomsection
\label{sga3:VIA:0.4}
设 \(G\) 是一个 \(A\)-概形.
若存在一个 \(A\)-态射
\[
  s:\operatorname{Spec}A\longrightarrow G
\]
将 \(\operatorname{Spec}A\) 的唯一点映到 \(G\) 的点 \(g\),
即同态 \(A\to\mathcal O_{G,g}\) 有一个收缩,
则称 \(g\) \emph{在 \(A\) 上严格有理 (strictly rational)}.
此时有 \(\kappa(g)=k\), 所以 \(g\) 是 \(G\) 的一个闭点.
事实上, 若 \(B\) 是 \(g\) 的某个仿射开邻域的环,
\(\mathfrak p\) 是 \(B\) 中对应于 \(g\) 的素理想,
则 \(B/\mathfrak p\subset k\) 是有限 \(A\)-代数,
从而是整的 Artin 环, 因而是域.

此后假设 \(G\) 是一个 \(A\)-群.
这样的态射 \(s:\operatorname{Spec}A\to G\)
定义了概形 \(G\) 在 \(A\) 上的一个自同构 \(r_s\),
称为\emph{由 \(s\) 给出的右平移}.
对每个态射 \(\pi:S\to\operatorname{Spec}A\),
\(G(S)\) 的自同构 \(r_s(\pi)\) 由下式定义:
\[
  x\longmapsto x\cdot G(\pi)(s),
  \qquad x\in G(S).
\]
同样地, \(\ell_s\) 表示\emph{由 \(s\) 给出的左平移},
即由下式定义的 \(G\) 的自同构:
\[
  \ell_s(\pi)(x)=G(\pi)(s)\cdot x,
  \qquad x\in G(S).
\]

概形 \(G\otimes_A k\) 与 \(G\) 有相同的底拓扑空间
\(\underline G\), \(G\otimes_A k\) 是一个 \(k\)-群,
且 \(s\otimes_A k\) 只依赖于 \(g\), 而不依赖于 \(s\).
因此, \(r_s\) 与 \(\ell_s\)
(或 \(r_{s\otimes k}\) 与 \(\ell_{s\otimes k}\))
所诱导的 \(\underline G\) 的自同构只依赖于 \(g\), 而不依赖于 \(s\).
所以, 若 \(P\) 是 \(\underline G\) 的子集, 可以记
\[
  r_g(P)\ \text{或}\ P\cdot g
  \qquad
  \bigl(\text{分别地, }\ell_g(P)\ \text{或}\ g\cdot P\bigr)
\]
来代替 \(r_s(P)\) (分别地, \(\ell_s(P)\)),
这与 \hyperref[sga3:VIA:0.1]{0.1} 一致.

\medskip
\noindent\textbf{注 0.4.1.}
\phantomsection
\label{sga3:VIA:0.4.1}
\textit{若 \(g\) 是 \(G\) 的一个严格有理点,
且 \(A\to A'\) 是 Artin 局部环态射,
则 \(G'=G\otimes_A A'\) 在 \(g\) 上方有唯一的点 \(g'\),
且 \(g'\) 在 \(A'\) 上严格有理.
此外, 若 \(P'\) 表示 \(P\) 在 \(G'\) 中的逆像,
则 \(P'\cdot g'\) 是 \(P\cdot g\) 的逆像 (参见 EGA I, 3.4.8).}

\medskip
\noindent\SGAInlineRefTarget{sga3:vi:proposition:0-5}\textbf{命题 0.5.}
\footnote{编者注: 原文在 \(G\) 于 \(A\) 上局部有限型的假设下陈述此结果.
由于此结果在一般情形中也有用, 且证明本质上相同,
此处在完全一般的情形下陈述并证明它. 下文将多次使用它.}
\phantomsection
\label{sga3:VIA:0.5}
\textit{设 \(G\) 是一个 \(A\)-群,
\(U,V\) 是 \(G\) 的两个稠密开子集.
则 \(U\cdot V\), 即 \(U\times_A V\) 在乘法下的像,
是 \(G\) 的整个底空间.}

\begin{proof}
由于 \(G\) 与 \(G\otimes_A k\) 有相同的底空间,
必要时以 \(k\) 代替 \(A\), 以 \(G\otimes_A k\) 代替 \(G\),
可以假设 \(A=k\).
设 \(g\in G\), 并置 \(K=\kappa(g)\).
则左平移 \(\ell_g\) 是 \(G_K\) 的自同构.
由于投影 \(G_K\to G\) 是开映射,
\(U_K\) 与 \(V_K\) 是 \(G_K\) 的稠密开子集,
\(V_K\) 在 \(\ell_g\) 下的像也如此.
因此存在 \(v\in V_K\), 使得 \(u=\ell_g(v)\) 属于 \(U_K\).

设 \(L\) 是包含 \(\kappa(v)\), 因而也包含 \(\kappa(u)\) 的
\(K\) 的一个扩域; 设 \(g_L\) 和 \(v_L\) 是由 \(g\) 和 \(v\)
诱导的 \(G_L\) 的 \(L\)-值点. 则
\[
  g_L\cdot v_L=u'
\]
是 \(G_L\) 中位于 \(u\) 上方的点, 因而
\[
  g_L=u'\cdot v_L^{-1}
\]
位于 \(U\cdot V\) 上方.
所以 \(g\in U\cdot V\), 命题得证.
\end{proof}

% Authority: Exp6A local p. 4; combined-reader p. 307; printed p. 296.
\medskip
\noindent\textbf{推论 0.5.1.}
\footnote{编者注: 此推论插入于此;
参见下文 \SGARefLink{prop:sga3-VIA-2.4}{2.4} 和
\SGARefLink{cor:sga3-VIA-2.6.2}{2.6.2}.}
\phantomsection
\label{sga3:VIA:0.5.1}
\textit{若 \(G\) 是不可约的 \(A\)-群, 则 \(G\) 拟紧.}

\begin{proof}
设 \(U\) 是 \(G\) 的一个非空仿射开子集.
则 \(U\) 在 \(G\) 中稠密, 所以由 \hyperref[sga3:VIA:0.5]{0.5}, 态射
\[
  \mu:U\times_A U\longrightarrow G
\]
是满射. 因此, 由于 \(U\times_A U\) 拟紧, \(G\) 也拟紧.
\end{proof}

\medskip
\noindent\textbf{推论 0.5.2.}\footnote{编者注: 此推论为新增内容,
在 VIII, \SGARefLink{sga3:viii:corollary:6:7}{6.7} 中隐含地使用了它;
另见 VI\textsubscript{B}, \SGARefLink{cor:sga3-VIB-6.2.5}{6.2.5}.}
\phantomsection
\label{sga3:VIA:0.5.2}
\textit{设 \(G\) 是一个 \(A\)-群概形,
\(H\) 是 \(G\) 在 \(A\) 上的一个群子概形. 则 \(H\) 是闭的.}

\begin{proof}
设 \(k'\) 是 \(A\) 的剩余域 \(k\) 的完美闭包.
由于 \(G\) 与 \(H\) 的底拓扑空间在基变换
\(A\to k\to k'\) 下不变, 可以假设 \(A=k\) 是完美域.
于是还可假设 \(G\) 与 \(H\) 约化, 因而几何约化.

设 \(\overline H\) 是 \(H\) 的闭包.
则 \(\mu^{-1}(\overline H)\) 是 \(G\times G\) 中包含
\(H\times H\) 的闭子集. 态射
\[
  H\longrightarrow\operatorname{Spec}k
  \qquad\text{及}\qquad
  \overline H\longrightarrow\operatorname{Spec}k
\]
都是普遍开映射.
由于 \(H\) 在 \(\overline H\) 中稠密,
\(H\times H\) 在 \(\overline H\times H\) 中稠密,
且 \(\overline H\times H\) 在
\(\overline H\times\overline H\) 中稠密.
所以 \(H\times H\) 在 \(\overline H\times\overline H\) 中稠密. 因而
\[
  \mu(\overline H\times\overline H)\subset\overline H.
\]
由于 \(\overline H\times\overline H\) 约化, \(\mu\) 诱导一个态射
\[
  \mu':\overline H\times\overline H\longrightarrow\overline H.
\]

现在设 \(g\in\overline H\), 并置 \(K=\kappa(g)\).
投影 \(\overline H_K\to\overline H\) 是开映射,
所以 \(H_K\) 与 \(\ell_g(H_K)\) 是
\(\overline H_K\) 的两个稠密开子集.
因此存在 \(u,v\in H_K\), 使得 \(\ell_g(v)=u\).
如 \SGARefLink{sga3:vi:proposition:0-5}{0.5} 的证明,
可得 \(g\) 属于 \(H\cdot H=H\). 因此 \(\overline H=H\).
\end{proof}

% Exact continuation cursor: Exp6A local p. 4, immediately before
% "1. Local properties of an A-group locally of finite type".

% Translation coverage: entire assigned component, retaining all source statements and proofs.
% Statistics: source lines=369; target lines=375; inline math=294; display math=18; unresolved translation TODOs=0; independent compilation not performed here.
